\documentclass[11pt]{article}
\usepackage[margin=0.68in]{geometry}
\usepackage{booktabs}
\usepackage{graphicx}
\usepackage{float}
\usepackage{newtxtext,newtxmath}
\usepackage{microtype}
\pagestyle{plain}
\begin{document}

\begin{table}[H]
\centering
\caption{\textbf{Original summary statistics for the underlying predictor
data.} Each row uses country-years in which both sector-specific predictors
are observed. Coverage is the percentage of the balanced 42-country period
that has a complete pair; debt and real asset-price growth are three-year
changes. The correlation and R-Zone share describe whether extreme debt and
price growth occur jointly. The main takeaway is that later data have broader
country coverage but substantially lower average debt growth and fewer joint
upper-tail observations than the historical estimation years. This change in
the predictor distribution helps explain why updated R-Zone frequencies and
crisis estimates need not equal their historical counterparts.}
\label{tab:data-overview}
\small
\input{../_output/data_overview_summary.tex}

\medskip
\begin{minipage}{0.97\linewidth}
\footnotesize
\textit{Notes:} $N$ is the number of complete country-year predictor pairs.
Coverage divides $N$ by 42 times the number of calendar years in the displayed
period. Debt growth is the three-year percentage-point change in debt/GDP;
price growth is 100 times the three-year log change in the relevant real asset
price. Corr. is the within-period Pearson correlation. R-Zone percentages use
the frozen historical Q80 debt and T66.7 price cutoffs.
\end{minipage}
\end{table}

\begin{figure}[H]
\centering
\includegraphics[width=0.86\textwidth]{../_output/data_overview_figure.jpg}
\caption{\textbf{Coverage and joint predictor distributions behind the
R-Zone classifications.} Panel A shows how many of the 42 countries have both
required predictors in each year; the shaded region begins after the common
historical forecast-origin window. Panels B and C plot every complete later
country-year pair and identify observations exceeding both frozen historical
thresholds. The reader should take away two features: coverage expands sharply
over time, so annual aggregates are based on changing country counts, and most
later observations lie outside the joint upper-right tail. R-Zones are
therefore genuinely rare combinations of rapid debt and asset-price growth,
not merely years in which either series is high on its own. Axes in the scatter
panels are limited to the first and 99th percentiles for readability; the table
uses all untrimmed observations.}
\label{fig:data-overview}
\end{figure}

\end{document}
